\chapter{Linguistic Categories, Online Processing, and Reasoning}
\label{sec:p19-language-categories}

Language can change what information a person receives, the distinctions
that must be expressed, and the resources available while solving a task.
These effects are not one hypothesis. A message can supply a new premise
without creating the capacity to entertain that premise. An utterance can
direct attention to a feature without making other features inaccessible.
An efficient verbal memory code can improve performance without replacing
the represented domain by a linguistic one.

Gleitman and Papafragou develop these distinctions in their author draft
of \emph{New Perspectives on Language and Thought}
\cite[author-draft pp.~3--17,56--61]{GleitmanPapafragouLanguageThoughtAuthorDraftLC}.
The draft is the presentation used here; its arguments are not silently
identified with the final Oxford chapter or their earlier Cambridge chapter.
The original experiments treated below provide separately attributable
evidence about colour discrimination, motion-event inspection, and the
construal of objects and substances. These are constrained empirical
questions, not tests of the logical validity of an inference.
The interpretation of a conjunction question is treated separately in
Chapter~\ref{sec:p19-conjunction-language}.
\index{language and thought}\index{linguistic relativity}
\index{linguistic determinism}\index{online processing}

\section{Specify what language is alleged to change}
\label{sec:p19-lc-claims}

\begin{hbdefinition}{def:p19-lc-evidence-record}{A language-effect evidence record}
Identify the linguistic contrast, the participants' language histories,
the task and its instructions, the stimulus information, the response
measure, and the processing interval to which the proposed effect applies.
Keep the production of a description, the interpretation of an instruction,
attention to a displayed event, immediate discrimination, and later memory
as different outcomes.

In the empirical reports below, \(N\) denotes a participant count, \(F\)
an analysis-of-variance test statistic, \(t\) a reported Student test
statistic, \(p\) its reported significance probability, and \(r\) a
reported Pearson correlation. Numbers in parentheses after \(F\) or \(t\)
are the reported degrees of freedom, not additional participants.
SEM is the reported standard error of a mean;
it is not an individual-response standard deviation or a confidence interval.
Condition means and item means are not independent observations merely
because a table has several rows.
\end{hbdefinition}

The strongest linguistic-determinist claim makes some conceptual content
unavailable without the relevant language or construction. A weaker
representational claim says that acquiring a language durably changes
category boundaries or the weighting of features, including when language
is not currently used. A processing claim instead concerns how existing
representations are selected, encoded, rehearsed, or interpreted in a task.
The latter can be important even if temporary: a short-lived change at the
time a decision is made may change its outcome.

The author draft does not treat contemporary positions as a simple division
between an initially empty mind and an unmodifiable mind. Its question is
how prior capacities and experience interact, and which causal contribution
is made by linguistic experience in particular. It also distinguishes
Whorf-inspired theoretical claims from an intellectual biography of Whorf;
his writings are not represented as one uniformly stated experimental
hypothesis \cite[author-draft pp.~5--10,17]{GleitmanPapafragouLanguageThoughtAuthorDraftLC}.

\begin{hbexample}{def:p19-lc-expression-countercase}{Failure to encode is not failure to represent}
A speaker says that a relative is an uncle. The English noun leaves open
whether the relationship is through the speaker's mother or father, and
through blood or marriage. The speaker may nevertheless know which
relationship holds. Conversely, one ambiguous word can be associated with
several distinguishable meanings. Thus a many-to-one linguistic expression
does not by itself establish a many-to-one conceptual representation.
This is an argument against identifying the two representations, not an
experiment showing that language never changes concepts.
\end{hbexample}

Pronouns, indexicals, omitted arguments, figurative reference, and repair
of misspeaking require an interpreter to combine linguistic information
with context and background knowledge. A description of rain, for example,
does not settle every spatial and temporal parameter independently of
context. The draft uses such cases to argue that linguistic encoding
underdetermines intended content. Whether a particular parameter belongs
in semantic structure or pragmatic enrichment remains a separate disputed
question; the argument does not resolve every semantics--pragmatics boundary
\cite[author-draft pp.~10--17]{GleitmanPapafragouLanguageThoughtAuthorDraftLC}.

\section{From a language contrast to a causal claim}
\label{sec:p19-lc-identification}

A cross-language comparison first establishes a linguistic difference and
then asks whether a concordant behavioural difference exists. Concordance
alone leaves several directions of explanation available. Linguistic
practice might influence a strategy; shared environmental practices might
influence both language and strategy; or a task might elicit different
interpretations from the two groups. Population comparisons also carry
differences in education, bilingualism, recruitment, and familiarity with
laboratory equipment.

Three further contrasts have different uses. Requiring a linguistic response
tests what happens when language is explicitly recruited. Interfering with
verbal processing tests whether an effect depends on resources occupied by
the secondary task. Comparing early and late behaviour can distinguish
event inspection from preparation to describe or remember it. None is a
pure intervention on an isolated faculty: a verbal task also imposes load,
and a change in instructions can change goals as well as wording.
The author's causal argument therefore depends on converging task-specific
contrasts, not on relabelling any observed group difference as an effect
of language \cite[author-draft pp.~57--60]{GleitmanPapafragouLanguageThoughtAuthorDraftLC}.

There are two particularly consequential processing routes. A listener may
use learned distributional regularities to interpret an underspecified
instruction or novel word. A viewer may also recode a nonverbal display
linguistically while preparing a response or preserving it in memory.
The second route can occur without overt speech. Consequently, a task
with pictures and key presses is not established as language-free merely
by its response format.

\section{Colour discrimination: the Russian-blue experiment}
\label{sec:p19-lc-colour}

Winawer and colleagues compared English's broad basic blue category with
the Russian distinction between \emph{goluboy} and \emph{siniy}.
Their task asked for a physically correct match, not a similarity rating
or a choice about how a colour ought to be named
\cite[pp.~7780--7785]{WinawerEtAl2007RussianBluesLC}.

\subsection{Participants, displays, and interventions}

The recruited sample contained 26 native Russian and 24 native English
speakers from the Boston area, tested at MIT. Russian participants had
begun acquiring English between ages 7 and 21. All received English
instructions. The paper reports mean ages of 28.9 and 26.3 years,
respectively. These are not randomized language groups.

Twenty displayed colours ranged from lighter to darker blues. On each
trial, one upper square and two lower squares remained visible until
response. One lower square was physically identical to the upper square.
The participant selected its side as quickly and accurately as possible.
Near distractors were two steps away on the stimulus continuum; far
distractors were four steps away. These steps are properties of the selected
stimulus set, not a claim of a universally uniform perceptual metric.

Each participant completed three 136-trial blocks, with order varied
randomly. One block had no secondary task. In the verbal block, an
eight-digit string was shown for three seconds and silently rehearsed
across eight colour trials; recall was then tested against a one-digit
foil. In the spatial block, participants retained a pattern of four shaded
cells in a four-by-four grid, then selected it against a foil differing in
one cell. The secondary tasks were pretested for similar recall accuracy;
that does not establish equivalence of every kind of resource demand.

After the discrimination blocks, each participant classified the twenty
colours twice, using the Russian words or English light/dark-blue labels.
The analysis located each person's boundary from these responses, with
classification latency used when the transition was ambiguous. The nine
near and nine far comparisons closest to that boundary were analysed.
Comparisons on opposite sides, or including the boundary stimulus, were
classified as between-category; two colours on the same side were within-category.
An English participant's light/dark-blue division was an instructed
division; it was not treated as an English basic-colour distinction.

Trials were removed for incorrect colour answers, failed secondary-task
answers, or colour-response latency above three seconds. Twelve percent
of trials were removed under these criteria. Participants losing at least
one quarter of trials were excluded: five Russian and three English
participants. The resulting analyses contained 21 participants per group.
The exclusion rule and the post-task classification are part of the
inferential procedure, not incidental details.

\subsection{What was measured}

Table~\ref{tab:p19-lc-blue-near} gives the reported near-colour means and
SEMs. A category advantage is the within-category mean latency minus
the between-category mean latency for the same condition; positive values
mean faster responses between categories. It is not a measure of how
many new colours a participant can see.

\begin{table}[htbp]
\centering
\begin{tabular}{lrrrr}
\hline
 &\multicolumn{2}{c}{Russian}&\multicolumn{2}{c}{English}\\
Condition&Between&Within&Between&Within\\ \hline
No secondary task&1164 (66)&1288 (77)&900 (51)&914 (52)\\
Spatial task&1162 (58)&1270 (56)&911 (41)&922 (46)\\
Verbal task&1325 (55)&1260 (50)&952 (41)&955 (46)\\ \hline
\end{tabular}
\caption{Near-colour matching latencies in milliseconds, with SEM in
parentheses; original Table 1 in Winawer and colleagues.}
\label{tab:p19-lc-blue-near}
\end{table}

The far-colour comparisons used the same conditions and participants.
Their complete reported means are in Table~\ref{tab:p19-lc-blue-far}.
The language-linked latency pattern was concentrated in the near-colour
comparisons, not a reliable between-category advantage at every distance.

\begin{table}[htbp]
\centering
\begin{tabular}{lrrrr}
\hline
 &\multicolumn{2}{c}{Russian}&\multicolumn{2}{c}{English}\\
Condition&Between&Within&Between&Within\\ \hline
No secondary task&998 (55)&999 (55)&758 (36)&735 (32)\\
Spatial task&1096 (64)&1096 (53)&819 (37)&835 (43)\\
Verbal task&1146 (62)&1132 (50)&831 (41)&821 (35)\\ \hline
\end{tabular}
\caption{Far-colour matching latencies in milliseconds, with SEM in
parentheses; the remaining conditions of the original Table 1.}
\label{tab:p19-lc-blue-far}
\end{table}

For Russian participants, the near-colour advantage without interference
was 124 milliseconds, with \(t(20)=2.59\), \(p=.0176\); the spatial
condition retained an advantage, whereas the verbal condition did not.
The reported category-by-interference interaction was
\(F(2,40)=5.3\), \(p=.009\). Its dependence on near rather than far
comparisons was reflected in the three-way interaction
\(F(2,40)=3.3\), \(p=.049\).
The negative verbal-condition difference is not established as a reversed
effect: its reported comparison had \(p=.076\).
English participants showed no detected category advantage in any block.
The paper also reports between-group contrasts in the changes caused by
interference, rather than relying solely on significance in one group and
non-significance in another
\cite[pp.~7782--7783,Table~1]{WinawerEtAl2007RussianBluesLC}.

Overall accuracy was high: the reported English and Russian means were
96.5\% and 95.7\%. A near-colour accuracy difference in the no-interference
Russian condition went in the opposite direction to the latency advantage.
The authors consequently examined speed--accuracy tradeoffs, and reported
that controlling for their accuracy contrast did not remove the association
between language group and the latency contrast. The result should not be
restated as simultaneous improvement on every accuracy and speed measure.

\subsection{The supported inference and its limits}

The display remained available, the physical matching question had a
correct answer, and verbal interference altered the relevant latency
pattern differently from the spatial control. This is evidence for an
online contribution of language-linked resources even in an objective
task with little retention demand. It therefore goes beyond showing that
people choose identically named chips when an experimenter asks an
ambiguous similarity question.

It does not show an inability of English speakers to distinguish these
blues: they performed the matching task accurately and supplied a similar
light/dark boundary. Nor did the experiment measure a permanent change in
visual processors. The original discussion explicitly leaves open whether
linguistic feedback affects earlier perceptual processing or whether a later
decision mechanism combines linguistic and sensory information. Both can
produce an online effect. Colour matching, remembered colour similarity,
and naming thus retain different evidential targets
\cite[pp.~7783--7785]{WinawerEtAl2007RussianBluesLC}.

\section{Motion events: inspection while preparing to speak}
\label{sec:p19-lc-motion-production}

English and Greek differ in their typical packaging of motion information.
In the selected studies, English tends to place manner in the main verb
and path elsewhere; Greek tends to encode bounded path in the verb and
manner in another expression. Neither language lacks the capacity to
express the other component. The differences are typological preferences
with construction-specific restrictions, not an absolute inventory of
available thoughts.

Papafragou, Hulbert, and Trueswell compared these packaging preferences
with the time course of eye movements during one experiment
\cite[pp.~155--184]{PapafragouHulbertTrueswell2008EventPerceptionLC}.
Seventeen native English and seventeen native Greek participants supplied
the analysed sample; three additional participants' records were discarded
for severe tracking loss. English participants were Pennsylvania psychology
undergraduates receiving course credit. Greek participants were students
or junior faculty in the Philadelphia area, paid for participation.
The authors note that the Greek group was fluent in English. Instructions
and introductory interaction were in the participant's native language.

Twelve test animations showed instrumental motion, such as a character
skating toward a stationary object. Each had a bounded and an unbounded
version; removing the goal produced the latter. Motion lasted three
seconds, followed by a beep and two seconds of the final still frame.
Each list contained six bounded and six unbounded test events among twelve
fillers, with counterpart and reverse-order lists. Participants were
randomly assigned to a list and either a description or a memorization
task. Description participants spoke after the beep. Memorization
participants knew that they would later recognize still images.

All received the final recognition test, unexpectedly for the description
group. Changed test images added or removed the goal; they did not
independently test fine changes in every manner and path feature.
Eye movements were recorded during viewing and subsequently coded
frame by frame. Manner regions were the spatially distinguishable
instruments; path regions were goal locations, including the corresponding
empty location in the unbounded version. Agent looks were coded separately.
These regions operationalize the selected event features. Looking at an
instrument is not a direct measurement of a complete manner concept, and
looking at an empty trajectory was not counted as path encoding.

\subsection{Temporal prioritization is not total attention}

Bounded-event descriptions used manner verbs in 72\% of English responses
and 36\% of Greek responses. In unbounded events the corresponding figures
were 76\% and 54\%. During preparation of bounded-event descriptions,
Greek participants prioritized the path endpoint earlier than English
participants, followed by a shift as the descriptions were assembled.
The first second already contained a language-by-region difference.
Unbounded events, where the linguistic contrast was less pronounced,
did not produce corresponding reliable gaze differences.

The global three-second totals did not show the same cross-language
contrast: what differed was the temporal order of inspecting components.
Both groups had a general path-endpoint preference. Thus the result must
not be described as English participants simply attending more to manner
over the entire event. Analyses of utterances mentioning manner early
also suggested that linear order of mention alone was insufficient to
explain the gaze pattern. This was an analysis of selected observed
utterances, not a randomized intervention on verb position
\cite[pp.~166--177,Table~1,Figs.~2--7]{PapafragouHulbertTrueswell2008EventPerceptionLC}.

In the memorization task, the groups' gaze patterns were similar while
motion unfolded. After the event stopped, they diverged: English
participants spent more time studying path endpoints. The authors
interpreted this later pattern as possible linguistic recoding of what
was less central in the ordinary verbal description. That interpretation
was a hypothesis about encoding, not a measured covert utterance.

Five participants were excluded from the recognition analysis for a strong
tendency to answer that changed pictures were unchanged. In the retained
sample, correct percentages for bounded and unbounded events were 92/91
for English and 92/92 for Greek in the description condition, and 90/88
and 73/71, respectively, in the memorization condition. These condition
means do not establish a universal memory deficit for Greek speakers.
They concern the added/removed-goal test, selected participants, and these
instructions. The later study below changed both the stimuli and the
recognition procedure.

The decisive comparison is therefore between production planning,
ongoing event inspection, and later encoding. A coarse end-of-task score
can conceal those stages, and a non-significant early contrast is not
proof that the underlying representations are identical in every respect.

\section{Motion events: load and availability of verbal coding}
\label{sec:p19-lc-motion-dual-task}

Trueswell and Papafragou conducted two further experiments designed to
test the proposed encoding route
\cite[pp.~64--82]{TrueswellPapafragou2010DualTaskLC}.
Experiment 1 recruited thirty English and thirty Greek university students,
approximately 18--22 years old. Greek participants were tested in Athens;
English participants were from Delaware or Pennsylvania. Participants
normally lived in the corresponding language environment; most had no
immersive second-language experience. Ten participants per language were
assigned to each of three viewing conditions.

Again, twelve test events involved a moving character, an instrument,
and a stationary goal, with twelve fillers. Each event had a three-second
animation and two-second still phase. In contrast with the 2008 experiment,
the later recognition test used animations, and test changes separately
replaced either the instrument or the path endpoint. It did not merely
add or remove a goal. There were no preview practice examples of changes.

The three conditions were viewing without an additional task, repeated
rhythm tapping with a pencil, and repeating aloud three two-digit numbers.
The secondary task continued through motion and still phases. The authors
selected the tapping and number tasks using pilot judgments of difficulty
and compared their effects on recognition. Similar overall impairment
does not prove process-pure or exactly equal loads.

Automatic scoring used moving regions for the instrument, endpoint, and
agent, with occlusion handled by counting the visible foreground region.
Trials with more than one third tracking loss were removed. Participants
losing five or more trials were excluded; this removed two Greek eye
records and no English records. About 3\% of trials were excluded among
retained participants. Recognition and eye-movement denominators therefore
must not be assumed to be the same merely because the recruited sample
contained sixty people.

\subsection{Experiment 1: a conditional rather than pervasive effect}

The principal gaze contrast was cumulative time at the path endpoint
minus cumulative time at the manner instrument. These are not complementary
proportions: participants also looked at agents and background regions.
Under tapping, English participants increasingly studied the path endpoint
relative to manner, while Greek participants changed in the opposite
direction. The divergence developed during the second second.
The language-by-task pattern was not reproduced by the concurrent verbal
task or ordinary no-interference viewing. The paper calls this a
\emph{reverse-Whorfian} pattern because attention was directed toward
what was usually outside, rather than inside, the main verb.

The analysis used crossed participant/item models of cumulative gaze,
with linear and quadratic time terms and specified random effects.
Separate one-second subject and item analyses also supported the
language-by-interference contrast in the second interval.
There was a small model-dependent curvature difference under no interference
which did not recur in the corresponding interval analyses. It was not
treated as robust evidence of a general language-specific inspection bias
\cite[pp.~71--75,Tables~1--2,Figs.~3--5]{TrueswellPapafragou2010DualTaskLC}.

Both secondary tasks impaired recognition relative to ordinary viewing.
The mixed-logit recognition analysis found effects of interference and
change type, but no detected language interaction. The most conspicuous
gaze difference therefore did not demonstrably improve the corresponding
recognition score. Optional recoding and successful recoding are different
claims. The authors also recognized a floor concern: interference-condition
recognition was not far above the binary task's chance level.

\subsection{Experiment 2: preparing before interference starts}

Ten further Greek and ten further English participants from the same
populations viewed the events without a secondary task during motion.
The number-repetition task began only at the beep, during the two-second
still frame. Participants knew this schedule in advance.
Gaze divergence developed around the onset of interference, over roughly
2.5--4 seconds, again in the direction of studying components less
central to the usual linguistic description. The pattern is consistent
with preparing a verbal code before that resource becomes occupied;
it does not directly observe the representation assumed to be prepared.

Recognition improved to about 70\% compared with the earlier interference
conditions. Its analysis again did not detect a language-by-change-type
interaction. The new condition thus addressed the earlier floor concern
without supplying evidence of the predicted language-congruent memory
asymmetry \cite[pp.~76--78,Table~3,Figs.~7--8]{TrueswellPapafragou2010DualTaskLC}.

\subsection{A genuine disagreement about generalization}

The 2010 authors accept that the Russian-blue result implicates online
linguistic resources, but dispute the inference to pervasive automatic
use of labels in perceptual tasks. They propose that a task containing
only shades of blue may encourage use of the relevant lexical contrast,
and suggest testing a broader stimulus set. That proposed manipulation
was not performed in their motion experiments. Their argument also
draws on lexical-decision and Stroop research, which supplies additional
hypotheses rather than new participants in the present studies.

The disagreement does not require denying the observed colour effect.
It concerns the scope and mechanism inferred from it. The colour study
shows an effect in its objective matching task; the motion studies show
that verbal encoding can be conditional on load, timing, and task goals.
Neither establishes a single mechanism operating in all perceptual domains
\cite[pp.~79--80]{TrueswellPapafragou2010DualTaskLC}.

\section{Objects, substances, and linguistic category learning}
\label{sec:p19-lc-entity-accounts}

One entity permits different relevant construals. A wooden table may be
counted as a table or described as wood. Breaking it apart may destroy
the table while preserving the material. This distinction makes different
properties relevant to classification and quantification. It is not
adequately described as a choice between seeing and not seeing the same
physical stimulus.

Li, Dunham, and Carey separate three accounts
\cite[pp.~488--492]{LiDunhamCarey2009SubstanceLC}.
The lexical-semantics account attributes group differences to what nouns
mean in count/mass versus classifier languages. A general feature-weighting
account says that habitual grammatical distinctions alter the salience
of properties whenever an entity is represented. The lexical-statistics
account instead predicts a difference while interpreting an ambiguous
noun: English learners can use the prevalence of count nouns as evidence
about an unfamiliar noun's intended reference. That effect can interact
with available information about shape and material without changing
task-neutral weights in every representation of the entity.

The authors reject the equation between lacking English count/mass syntax
and lacking individuation. Classifier languages have devices for referring
to and quantifying individuals. Mandarin nouns for a fork and a spoon
need not denote one common metal essence simply because numerals require
classifiers. Conversely, English mass morphology is not invariably a
substance label. Grammar, lexical reference, and conceptual availability
must each be investigated.

\section{The six entity-construal experiments}
\label{sec:p19-lc-entity-experiments}

\subsection{Experiment 1: language groups and two instruction forms}

Thirty-two English-speaking Harvard students and thirty-two Mandarin-speaking
National Taiwan University students participated. Sixteen per group were
randomly allocated to a novel-label task and sixteen to a no-novel-label
task. There were twelve standards: four complex solids, four simple solids,
and four non-solids. Each was compared with a same-shape/different-material
alternative and a same-material/different-shape alternative. The choices
were between specified displayed alternatives, not unrestricted accounts
of everything the participant believed about the object.

The novel-label version introduced a new noun in a frame compatible with
English count and mass nouns. The other version pointed out a standard and
asked for the same alternative, without providing a new noun. Both required
interpreting instructions; the latter was not a demonstrated language-free
condition. Two practice triads required unambiguous shape and material
matches, and all participants answered them correctly.

English participants selected shape more often overall: 58.6\% versus
31.3\%. Complex solids, simple solids, and non-solids respectively
elicited English means of 94.5, 59.4, and 21.9\%, and Mandarin means
of 73.4, 16.4, and 3.6\%. There was no detected instruction effect.
The largest group difference occurred for the simple solids. The common
ordering of stimulus types coexisted with group differences in the amount
of shape-based selection
\cite[pp.~492--496,Table~1,Fig.~1]{LiDunhamCarey2009SubstanceLC}.

\subsection{Experiment 2: separating stimulus properties}

Sixteen English, sixteen Mandarin, and twenty Japanese adults supplied
the next sample. English and Mandarin participants were university
students; Japanese participants were recent graduates in their early
twenties recruited in Kirishima. No participant had taken part previously.

Sixteen handcrafted unfamiliar standards replaced the partly recognizable
items. Twelve solids crossed rated shape complexity with the rated
likelihood of a shape-dependent function; four were non-solid.
Each still had a shape and a material comparison. These entities had no
actual assigned function. Ratings concerned an observer's assessment of
possible functional relevance, not measured usefulness. The selected
properties were therefore partly conceptual rather than mere outline
complexity. All participants performed the novel-label task.

Among solids the reported shape-choice means were 65\% for English,
35\% for Mandarin, and 47\% for Japanese. Potential functional relevance
and shape complexity both influenced choices. Seeing two unusual copies
of a shape may itself provide evidence that its form is nonaccidental.
The design cannot treat shape complexity, multiple copies, and evidence
about function as interchangeable features
\cite[pp.~496--501,Table~2,Fig.~2]{LiDunhamCarey2009SubstanceLC}.

\subsection{Experiment 3: explicit construal ratings}

Sixteen English, sixteen Mandarin, and twenty-four Japanese adults,
again not previously tested, rated the same sixteen standards one at a
time. The procedure explained the object/substance contrast using a table
and ice cream, and made the importance or irrelevance of form explicit.
Participants used a seven-point scale. The administered scale ran from
object at one to substance at seven; the published analyses inverted it,
so higher analysed numbers indicate more object-like construals.
The scale's direction must not be inferred from an isolated graph axis.

For the solid standards, reported analysed means were 4.9, 5.1, and 4.7
for English, Mandarin, and Japanese, respectively. English participants
were not systematically more object-oriented here. Functional relevance
remained influential. Average item ratings correlated strongly across
groups: English--Mandarin \(r=.88\), English--Japanese \(r=.97\), and
Mandarin--Japanese \(r=.92\). A functionality-by-language interaction
and some Mandarin--Japanese differences remained; the data must not be
described as proving numerical equality of all ratings.

This contrast limits the universal-scope feature-weighting interpretation.
The rating task did not reproduce the earlier English shape advantage.
However, the task was not devoid of language: it supplied words and explicit
examples for the relevant conceptual distinction. It also changed the
response scale and removed comparison copies. It supports task dependence,
not a clean estimate of the causal effect of one deleted noun
\cite[pp.~501--505,517--518,Appendix~B]{LiDunhamCarey2009SubstanceLC}.

\subsection{Experiment 4: younger children}

The sample contained 192 children: thirty-two at each of ages two, three,
and four in each language group. English learners came from greater-Boston
preschools, daycares, and the laboratory database; Mandarin learners came
chiefly from Taiwan preschools. The reported socioeconomic range was middle
to upper-middle class. Children were individually tested and randomly
assigned within each age/language group to the two instruction forms.
They used the first experiment's twelve triads, with practice repeated
when necessary.

Children distinguished the stimulus types but did not exhibit the adult
group difference: overall shape selection was about 51\% for English
and 49\% for Mandarin learners. Across both languages, shape selection
for non-solids declined from 33\% at age two to 15\% and 11\% at
ages three and four. This is a cross-sectional pattern, not a longitudinal
trajectory of the same children. A common developmental change in
material sensitivity can coexist with a later language-related difference
\cite[pp.~505--509,Fig.~6]{LiDunhamCarey2009SubstanceLC}.

\subsection{Experiment 5: ages three through five}

The handcrafted sixteen-triad set was used with 58 English learners
(18, 19, 21 across ages three, four, five) and 52 Mandarin learners
(16, 20, 16). All performed the novel-label task. The oldest groups gave
more language-differentiated choices, while functional relevance influenced
both populations. Item-level comparisons showed a difference at ages four
and five, but the age interaction did not have the same significance in
subject and item analyses. A further split of four-year-olds associated
the difference with the older subgroup.

These details prevent converting the paper's proposed emergence around
late four to five into a sharp universal age threshold. It was a pattern
under the selected stimuli and analyses, with different onset estimates
in earlier studies. English and Mandarin children continued to show
strongly correlated item preferences
\cite[pp.~509--514,Figs.~7--8]{LiDunhamCarey2009SubstanceLC}.

\subsection{Experiment 6: Japanese children}

Fourteen Japanese four-year-olds and sixteen five-year-olds from a
Kirishima preschool performed the same task. Their choices were compared
with the already collected English age groups; the English children
were not a new sample of the sixth experiment. An item-level group
difference was detected among five-year-olds, but not four-year-olds.
Several subject-level main or interaction tests did not detect the same
differences. Functional relevance influenced both groups, and child
responses did not reproduce the adults' complexity effect
\cite[pp.~514--516,Figs.~9--10]{LiDunhamCarey2009SubstanceLC}.

\subsection{The explanatory conclusion remains conditional}

The paper supports a distinction between properties influencing construal
across languages and language-related shifts in the selected naming tasks.
The lexical-statistics interpretation combines learned expectations about
nouns with information about a referent's form. Its developmental explanation
adds increasing knowledge of materials, which could offset or interact with
the English noun bias. That two-source explanation is proposed, not itself
established by randomly manipulating all relevant learning histories.

The authors retain another alternative for similarity tasks: the contrast
may have broader scope in comparisons among entities even if it disappears
in explicit ratings. Their suggested change to the demonstrative wording
would test that possibility. It is not an already completed seventh
experiment. Neither the rating result nor the developmental results rule
out strong or weak language effects in other domains
\cite[pp.~517--520]{LiDunhamCarey2009SubstanceLC}.

\section{The author account across other domains}
\label{sec:p19-lc-domains}

Gleitman and Papafragou's author presentation relates these experiments
to several further research questions. These connections have different
standing from the original experimental records just specified.

\paragraph{Speech perception.}
The draft treats experience-dependent phonetic and phonological
organization as a substantive effect of linguistic learning, not something
that can be dismissed merely for occurring within language. Its examples
include sensitivity to exposure-language sound distinctions and the
interaction of morphology with the perceived sounds in words. It also
qualifies claims of permanent loss by later learning and training.
Such language-on-language learning does not by itself establish that
nonlinguistic categories must be reorganized in the same way
\cite[author-draft pp.~17--20]{GleitmanPapafragouLanguageThoughtAuthorDraftLC}.

\paragraph{Spatial relations and reference frames.}
The draft distinguishes differences in lexical classification of support,
containment, and tightness of fit from an inability to represent those
relations. It also questions rotation tasks whose instruction asks a
participant to reconstruct an array in an unspecified same way: preserving
world direction and preserving body-relative direction are different
answers. Landmark manipulations and unambiguous instructions bear on
interpretation and strategy rather than merely on vocabulary size.
Its cited Korean/English and Mayan/European experiments retain their
separate participants and designs; they are not replications performed by
the colour or motion samples
\cite[author-draft pp.~30--33,37--41]{GleitmanPapafragouLanguageThoughtAuthorDraftLC}.
The geometrical and landmark contributions to spatial performance are
developed in Chapter~\ref{sec:p19-spatial-orientation-and-updating}.

\paragraph{Evidentiality.}
The author account contrasts grammatical marking of information source
with knowing whether information was seen, inferred, or supplied by
another person. It discusses English/Korean child comparisons and a
possible developmental precedence of source understanding over mastery
of evidential morphology. This is a claim about the cited developmental
tasks, not proof that grammatical source markers never affect memory or
communication. A grammatical obligation to mark a source also does not
validate the source or confer knowledge on the speaker
\cite[author-draft pp.~41--42]{GleitmanPapafragouLanguageThoughtAuthorDraftLC}.

\paragraph{Temporal metaphor.}
The draft separates habitual spatial expressions for time from the
representations used in a primed temporal judgment. It discusses
cross-language priming, short training that changed the reported pattern,
and failed replication or corpus challenges to the original account.
Its point is that rapid reversibility constrains an interpretation in
terms of immutable conceptual limits. The disputed original results
remain disputed; training-induced flexibility does not establish that
every language-related temporal effect is absent
\cite[author-draft pp.~43--44]{GleitmanPapafragouLanguageThoughtAuthorDraftLC}.

\paragraph{Number and cross-domain combinations.}
The author account distinguishes exact matching while sets remain available
from exact retention across delay, and approximate quantity from exact
large-number calculation. It discusses linguistic count lists as possible
tools for encoding and communicating quantity and separately evaluates
the stronger claim that grammar creates the number system's generative
structure. A linguistic list's surface regularity, its successor relation,
and the procedure for calculating are not identical. The draft also
examines proposals that language is necessary for combining geometric
and object information. Such a proposal owes an account of how the
combined meaning is constructed; calling language a medium does not
by itself supply that account
\cite[author-draft pp.~44--56]{GleitmanPapafragouLanguageThoughtAuthorDraftLC}.

The draft ends with language as a means of selecting, organizing, and
retaining information, including labels that invite category rather than
thematic matches. Its preferred account is not the denial of all influence
but an account of when linguistic representations participate in processing.
The distinction between a communicated category, a useful label, and a
necessary condition on conceptual capacity remains essential.

\section{Using these results in the analysis of reasoning}
\label{sec:p19-lc-use}

\begin{hbmethod}{def:p19-lc-method}{Appraising a language-based explanation}
State a prediction that distinguishes capacity, durable representation,
and currently recruited processing. If a proposed conceptual distinction
is claimed to be unavailable, provide a task in which participants would
need that distinction, with unambiguous instructions and an independently
specified success criterion. A preferred response in an ambiguous task
does not establish incapacity to produce the alternative.

For a processing claim, specify the phase: understanding instructions,
assembling a spoken description, inspecting a stimulus, encoding it,
retaining it, or selecting a response. Compare an appropriately selected
nonverbal secondary task, and retain the secondary-task scores and
exclusion rules. An interference interaction supports dependence on the
affected resources; identifying the precise internal mechanism requires
additional discrimination among alternatives.

Record bilingual experience, recruitment, educational and cultural setting,
age, stimulus range, and response format. Preserve the distinction between
participant-level and item-level analyses and between randomized task
allocation and nonrandomized language membership. Do not turn a failure
to detect a difference into evidence of exact equivalence.

Use the finding at its measured scope. A colour-matching latency contrast
is not a demonstration that one population reasons more rationally; a
production preference is not a conceptual prohibition; and a memory aid
does not change the validity conditions of the reasoning performed with
the remembered information. Conversely, an effect's temporary character
does not make it irrelevant when that processing episode determines the
answer.
\end{hbmethod}
